% TOP_PROBLEM: {"id": 2927, "title": "Kirby Problem 4.51", "category_id": 7, "category": {"id": 7, "name": "topology", "display_name": "Topology", "description": "Properties preserved under continuous deformations.", "slug": "topology", "order_index": 7, "created_at": "2026-07-31T15:26:25.671Z"}, "status": "open"}
% FIRST_POSTED: null
% ENTRY_KIND: statement_only
% Imported from: batch13.tex; UnsolvedMath ID 2927
\documentclass[11pt]{article}
\usepackage[margin=25mm]{geometry}
\usepackage{fontspec}
\setmainfont{Latin Modern Roman}
\usepackage{amsmath,amssymb,mathtools}
\usepackage{unicode-math}
\setmathfont{Latin Modern Math}
\usepackage{xurl,hyperref}
\hypersetup{colorlinks=true,urlcolor=blue}
\setcounter{secnumdepth}{0}
\setlength{\parindent}{0pt}
\setlength{\parskip}{5pt}
\setlength{\emergencystretch}{3em}
\newcommand{\sref}[2]{\hyperlink{source:#1:#2}{[S#2]}}
\newcommand{\cataloguescope}[1]{\paragraph{Recorded target.}#1}
\begin{document}
% UnsolvedMath ID 2927; source catalogue code KP-4.51
\setcounter{section}{2926}
\section{Kirby Problem 4.51}\label{2927}
{\small\sffamily UnsolvedMath ID 2927 · Topology}
\subsection{Definitions and mathematical statement}

\paragraph{Shared notation.}$\mathbb N=\{1,2,\ldots\}$, and $\mathbb Z,\mathbb Q,\mathbb R,\mathbb C$ denote the integers, rationals, reals and complex numbers. Any other conventions are specified locally.


Let $X$ be a closed smooth four-manifold with $\pi_1(X)\cong\mathbb Z$, choose a generator $t$, and put $\Lambda=\mathbb Z[t,t^{-1}]$. Its universal cover $\widetilde X$ is oriented. Let $w(t)\in\{1,-1\}$ record whether the deck transformation preserves that orientation; the group-ring involution is
\[
\overline{t}=w(t)t^{-1}.
\]
Via the Hurewicz identification $\pi_2(X)\cong H_2(\widetilde X;\mathbb Z)$, deck transformations make this a $\Lambda$-module. For compact representatives of classes $a,b$, the equivariant intersection form is
\[
\lambda_X(a,b)=\sum_{j\in\mathbb Z}w(t)^j\,(a\cdot t^jb)\,t^j\in\Lambda.
\]
Only finitely many intersection numbers are nonzero. This is Hermitian: it is linear in the first argument, conjugate-linear in the second, and $\lambda_X(b,a)=\overline{\lambda_X(a,b)}$. For orientable $X$, $w(t)=1$ and the involution is the usual $t\mapsto t^{-1}$.

A form is extended from $\mathbb Z$ if it has a $\Lambda$-basis with integer Gram-matrix entries. Equivalently it is obtained by extending scalars from a symmetric integral form on a finite-rank free abelian group, using the displayed involution.
\paragraph{Statement recorded in the supplied source.}
For every closed smooth four-manifold $X$ with $\pi_1(X)\cong\mathbb Z$, is $\lambda_X$ extended from $\mathbb Z$? Precisely, must there be a $\Lambda$-basis $e_1,\ldots,e_r$ of $\pi_2(X)$ such that
\[
\lambda_X(e_i,e_j)\in\mathbb Z\qquad(1\le i,j\le r)?
\]
The target is an integral matrix for the full equivariant form, not only for its augmentation at $t=1$. \sref{2927}{2}\sref{2927}{3}
\subsection{Short English statement}
Must the full infinite-cyclic equivariant intersection form of every closed smooth four-manifold have an integral Gram matrix after changing its group-ring basis?
\subsection{Sources}
\begin{description}
\item[\hypertarget{source:2927:1}{S1}] ULAM.AI and the UnsolvedMath Contributors, UnsolvedMath: A Curated Collection of Open Mathematics Problems, record 2927 (KP-4.51). CC BY 4.0. \url{https://huggingface.co/datasets/ulamai/UnsolvedMath}.
\item[\hypertarget{source:2927:2}{S2}] R. İnanç Baykur, Robion C. Kirby and Daniel Ruberman, K3—A New Problem List in Low-Dimensional Topology, Mathematical Surveys and Monographs 295 (2026), author version, Problem 4.51 and Remarks, PDF page 230. \url{https://math.berkeley.edu/sites/default/files/surv-295-ruberman-watermarked-author-pdf.pdf}.
\item[\hypertarget{source:2927:3}{S3}] Stefan Friedl, Ian Hambleton, Paul Melvin and Peter Teichner, Non-smoothable four-manifolds with infinite cyclic fundamental group, arXiv:math/0611077v2 (2007), Introduction, Conjecture 1.3, page 2. \url{https://arxiv.org/pdf/math/0611077v2}.
\end{description}
\label{end:2927}
\end{document}
